\documentclass[../../../include/open-logic-section]{subfiles}

\begin{document}

\olfileid{sth}{spine}{Valphabasic}
\olsection{د پړاوونو بنسټيز خاصيتونه}

د $V_\alpha$ ګانو د تعريف د بنسټيز اهميت د څرګندولو دپاره به ځينې
بنسټيزې نتيجې وړاندې کړو. په يوه تعريف پېل کوو:\footnote{د دې خاصيت
دپاره معياري اصطلاح نشته؛ اصلي متن ئې «potent» بولي، د \citet{ButtonLT1}
د نومولو په تعقيب. دلته ئې پښتو بڼه «توانمن» اخيستل شوې ده.}
\begin{defn}
	سټ $A$ \emph{توانمن} دے، که او يوازې که $\forall x((\exists y \in A)x \subseteq y \lif x \in A)$ وي.
\end{defn}
\begin{lem}\ollabel{Valphabasicprops}
د هر ترتيبي عدد $\alpha$ دپاره:
\begin{enumerate}
	\item\ollabel{Valphatrans} هر $V_\alpha$ متعدي دے.
	\item\ollabel{Valphapotent} هر $V_\alpha$ توانمن دے.
	\item\ollabel{Valphacum} که $\gamma \in \alpha$ وي، نو $V_\gamma
	\in V_\alpha$ دے (او له دې $V_\gamma \subseteq V_\alpha$ هم د
	\olref{Valphatrans} له مخې دے).
\end{enumerate}
\end{lem}

\begin{proof}
دا د نامتناهيت هاخوا په هممهاله استقرا ثابتوو. د استقرا دپاره فرض کړئ
چې \olref{Valphatrans}--\olref{Valphacum} د هر ترتيبي عدد $\beta < \alpha$
دپاره صدق کوي.

د $\alpha = \emptyset$ حالت څرګند دے.

فرض کړئ $\alpha = \ordsucc{\beta}$ دے. د \olref{Valphacum} د ښودلو
دپاره، که $\gamma \in \alpha$ وي، نو $V_\gamma \subseteq V_\beta$ د فرض
له مخې دے، نو $V_\gamma \in \Pow{V_\beta} = V_\alpha$ دے.
د \olref{Valphapotent} د ښودلو دپاره فرض کړئ $A \subseteq B \in V_\alpha$
دي، يعنې $A
\subseteq B \subseteq V_\beta$؛ نو $A \subseteq V_\beta$ دے، نو $A \in
V_\alpha$ دے. د \olref{Valphatrans} د ښودلو دپاره وګورئ چې، که $x \in A \in
V_\alpha$ وي، نو $A \subseteq V_\beta$ لرو، نو $x \in V_\beta$ دے، نو
$x
\subseteq V_\beta$ دے، ځکه چې $V_\beta$ د فرض له مخې متعدي دے، نو $x
\in V_\alpha$ دے.

فرض کړئ $\alpha$ حدي ترتيبي عدد دے. د \olref{Valphacum} د ښودلو
دپاره، که $\gamma \in \alpha$ وي، نو $\gamma \in \ordsucc{\gamma} \in \alpha$
دي، نو $V_\gamma \in V_{\ordsucc{\gamma}}$ د فرض له مخې دے، او له دې
$V_\gamma \in \bigcup_{\beta \in \alpha} V_\beta = V_\alpha$ دے.
د \olref{Valphatrans} او \olref{Valphapotent} د ښودلو دپاره يوازې وګورئ
چې د متعدي سټونو اتحاد متعدي دے، او د توانمنو سټونو اتحاد توانمن دے.
\end{proof}

\begin{lem}\ollabel{Valphanotref}
د هر ترتيبي عدد $\alpha$ دپاره، $V_\alpha \notin V_\alpha$ دے.
\end{lem}

\begin{proof}
د نامتناهيت هاخوا استقرا کاروو. ښکاره ده چې $V_\emptyset \notin V_\emptyset$ دے.

که $V_{\ordsucc{\alpha}} \in V_{\ordsucc{\alpha}} = \Pow{V_\alpha}$ وي،
نو $V_{\ordsucc{\alpha}} \subseteq V_\alpha$ دے؛ او ځکه چې $V_\alpha
\in V_{\ordsucc{\alpha}}$ د \olref{Valphabasicprops} له مخې دے،
$V_\alpha \in V_\alpha$ لرو. په معکوسه توګه: که $V_\alpha \notin V_\alpha$
وي، نو $V_{\ordsucc{\alpha}} \notin V_{\ordsucc{\alpha}}$ دے.

که $\alpha$ حدي عدد وي او $V_\alpha \in V_\alpha = \bigcup_{\beta \in
\alpha}V_\beta$ وي، نو $V_\alpha \in V_\beta$ د کوم $\beta \in
\alpha$ دپاره دے؛ خو بيا $V_\beta \in V_\alpha$ هم دے، نو $V_\beta \in
V_\beta$ د \olref{Valphabasicprops} په دوه ځله کارولو دے.
په معکوسه توګه، که $V_\beta
\notin V_\beta$ د ټولو $\beta \in \alpha$ دپاره وي، نو $V_\alpha \notin
V_\alpha$ دے.
\end{proof}

\begin{cor}
د هر دوو ترتيبي عددونو $\alpha, \beta$ دپاره: $\alpha \in \beta$، که او يوازې که $V_\alpha \in V_\beta$ وي.
\end{cor}

\begin{proof}
\Olref{Valphabasicprops} يو لورے ورکوي. په معکوسه توګه فرض کړئ
$V_\alpha \in V_\beta$ دے. نو $\alpha \neq \beta$ د \olref{Valphanotref}
له مخې دي؛ او $\beta \notin \alpha$، ځکه چې که داسې نۀ وي، نو
$V_\beta \in V_\alpha$ به لرو او له دې $V_\beta \in V_\beta$ د
\olref{Valphabasicprops} په دوه ځله کارولو لرو، چې له
\olref{Valphanotref} سره په تناقض کښې دي. نو $\alpha \in \beta$ د درې ګونتوب له مخې دے.
\end{proof}
\noindent
دا ټول موږ ته اجازه راکوي چې هر $V_\alpha$ د مراتبو د لړۍ د $\alpha$
پړاو په توګه وګڼو. وجه ئې دا ده:

بې شکه $V_\alpha$ ګانې د يوه \emph{تکراري} بهير په وسيله جوړې شوې
ګڼلے شو، ځکه چې د ترتيبي عددونو کارول مو د تکرار تصور تعقيبوي.
همدارنګه، که يو پړاو تر بل مخکې جوړېږي، يعنې
$V_\beta \in V_\alpha$، يعنې $\beta \in \alpha$، نو د جوړولو بهير مو
\emph{تجمعي} دے، ځکه چې $V_\beta \subseteq V_\alpha$ دے.
په پای کښې، په رښتيا د هغو سټونو \emph{ټولې} ممکنې ټولګې جوړوو چې
په کوم مخکيني پړاو کښې موجود وو، ځکه چې هر تالي پړاو
$V_{\ordsucc{\alpha}}$ د خپل مخکيني پړاو $V_\alpha$ د فرعي سټونو سټ دے.

په دې توګه، په $\ZFminus$ کښې د سټونو د تجمعي--تکراري مراتبو د لړۍ
د تصور په کره بيانولو کښې \emph{تقريباً} کار بشپړ شوے دے.
خو، بې شکه، لا د تعويض توجيه کول پکار دي.

\end{document}
